\documentclass[12pt, letterpaper]{article}

%-------Packages---------
\usepackage{amssymb,amsfonts}
\usepackage{mathptmx}
\usepackage{amsthm}
\usepackage{graphicx}
\usepackage[all,arc]{xy}
\usepackage[colorlinks=true, allcolors=blue]{hyperref}
\usepackage{enumerate}
\usepackage{bbm}
\usepackage{dsfont}
\usepackage{mathrsfs}
\usepackage{tikz-cd}
\usepackage{setspace}
\usepackage{import}
\usepackage[utf8]{inputenc}
\usepackage{hyperref}
\usepackage{tikz}
\usepackage{float}
\usepackage{mdframed}
\usepackage{fancyhdr}
\usepackage[nointegrals]{wasysym}

\usepackage[left=1in,top=1in,right=1in,bottom=1in]{geometry}

\usepackage{mathtools}
\usepackage{setspace}
\usepackage{cleveref}
\usepackage{multicol}

%--------Theorem Environments--------
%theoremstyle{plain} --- default
\newmdtheoremenv{theorem}{Theorem}[subsection]
\newmdtheoremenv{corollary}[theorem]{Corollary}
\newtheorem{proposition}[theorem]{Proposition}
\newmdtheoremenv{lemma}[theorem]{Lemma}
\newtheorem{conj}[theorem]{Conjecture}
\newtheorem{quest}[theorem]{Question}

\theoremstyle{definition}
\newmdtheoremenv{definition}[theorem]{Definition}
\newmdtheoremenv{definitions}[theorem]{Definitions}
\newtheorem{example}[theorem]{Example}
\newtheorem{algorithm}[theorem]{Algorithm}
\newtheorem{rmk}[theorem]{Remark}
\newtheorem{exercise}[theorem]{Exercise}

%----------- my commands -------------
\newcommand{\C}{\mathbb{C}}
\newcommand{\R}{\mathbb{R}}
\newcommand{\Q}{\mathbb{Q}}
\newcommand{\Z}{\mathbb{Z}}
\newcommand{\N}{\mathbb{N}}
\renewcommand{\P}{\mathbb{P}}
\newcommand{\contr}{\Rightarrow \Leftarrow}
\newcommand{\HRule}[1]{\rule{\linewidth}{#1}}
\newcommand{\s}{\(\,\)}
\renewcommand{\t}[1]{\text{#1}}
\newcommand{\ang}[1]{\left\langle #1 \right\rangle}

% Meta data
\setstretch{1.2}

\begin{document}

\pagestyle{fancy}
\fancyhf{}
\fancyhead[LOE]{\slshape{\textbf{Vincent Ferrara}}}
\fancyhead[ROE]{\thepage}

\subsection*{Problem 2}
\[
\P(1 \text{ wins on first round})=\dfrac{6}{36}=\dfrac{1}{6}
\]
\[
\P(2 \text{ wins on first round})=\dfrac{1}{36}+\dfrac{2}{36}+\dfrac{3}{36}=\dfrac{1}{6}
\]
\[
\P(\text{tie on first round})=\P(1 \text{ loses on first round})\P(2 \text{ loses on first round})=\left( \dfrac{5}{6} \right)^2=\dfrac{25}{36}
\]
\[
\P(1 \text{wins})=\P(1 \text{ wins round})+\P(\text{tie})\P(1 \text{ wins round})+\P(\text{ties} 2)\P(1 \text{ wins round})+\cdots
\]
\[
\P(1 \text{wins})=\dfrac{1}{6}+\dfrac{1}{6}\dfrac{25}{36}+\dfrac{1}{6}\left( \dfrac{25}{36} \right)^2+\cdots
\]
\[
\P(1 \text{wins})=\dfrac{\frac{1}{6}}{1-\frac{25}{36}}=\dfrac{6}{11}
\]

\newpage
\subsection*{Problem 6}
\begin{enumerate}
    \item There is a total of $13!$ permutations to order 13 cards.\\
    To count the ways that the Jack, King, and Queen are next to each other we can treat them as 1 card,
    so there is $3!$ ways to order those cards and $11!$ ways to order the 10 cards plus our group of 3 thus
    \[
    \P(E)=\dfrac{3! \times 11!}{13!}
    \]
    \item Again there is a totall of $13!$ permutations to order 13 cards.\\
    There are 3 ways to pair each element KQ, KJ, QJ
    and using a proof similar to in 1. we know\\
    \[
    \P(KQ)=\dfrac{2! \times 12!}{13!}, \quad
    \P(KJ)=\dfrac{2! \times 12!}{13!}, \quad
    \P(QJ)=\dfrac{2! \times 12!}{13!}
    \]
    $\P(KQ \cap KJ) \implies$ we either have $QKJ$ or $JKQ$ thus\\
    \[
    \P(KQ \cap KJ)=\dfrac{2! \times 11!}{13!}, \quad
    \P(KQ \cap QJ)=\dfrac{2! \times 11!}{13!}, \quad
    \P(QJ \cap KJ)=\dfrac{2! \times 11!}{13!}
    \]
    $\P(KQ \cap KJ \cap QJ)=0$ since we can't make them all adjacent to each other.\\
    Thus, by inclusion exclusion\\
    \[
    \P(E)=1-\P(KQ)-\P(KJ)-\P(QJ)+\P(KQ \cap KJ)+\P(KQ \cap QJ)+\P(QJ \cap KJ)-\P(KQ \cap KJ \cap QJ)
    \]
    \[
    \P(E)=1-3\dfrac{2! \times 12!}{13!}+3\dfrac{2! \times 11!}{13!}
    \]
    \item Again there is a total of $13!$ permutations to order 13 cards.\\
    The number of ways to count this is first we start with $13!$ but since there is $3!$ ways to order the K,Q,J and we only care about one
    we over count by a factor of $3!$ thus the number of ways is $13!/3!$ thus $\P(E)=\dfrac{\frac{13!}{3!}}{13!}=\dfrac{1}{6}$
\end{enumerate}

\newpage
\subsection*{Problem 7}
\begin{enumerate}
    \item LHS: This is the number of ways to pick $n$ objects from $2n$ objects\\
    RHS: We first split our $2n$ objects into sets of $n$ objects we then take $k$ objects from the first $\binom{n}{k}$ and take $n-k$ from the second $\binom{n}{k}=\binom{n}{n-k}$
    we then sum from 0 to $n$ to get all combinations of picking from set 1 and set 2 thus the sum is equivalent to counting the number of ways to pick $n$ objects from $2n$ objects
    \item LHS: This is the number of ways to make $m$ long words with $3$ letter alphabet lets say $a,b,c$
    RHS: We first set all letters to $a$ we then pick $k$ letters from our $m$ $\binom{m}{k}$ then we pick either $b$ or $c$ for those $k$ letters $2^k$.
    We then sum from 0 to $m$ to get all combinations of picking from the letter $a$'s thus the sum is equivalent to counting the number of ways to make $m$ long words with 3 letter alphabet
\end{enumerate}

\newpage
\subsection*{Problem 9}
\begin{enumerate}
    \item $\P(\text{first head on  the $n$-th throw})=\P(n-1 \t{ tails})\P(\t{heads})=\left( \dfrac{1}{2} \right)^{n-1}\times \dfrac{1}{2}=\left( \dfrac{1}{2} \right)^n$
    \item $\P(\# \t{ of heads and tails are equal on $n$-th throw})$\\
    Total number of posibilities is $2^n$\\
    If $n$ is odd $\P(\# \t{ of heads and tails are equal on $n$-th throw})=0$\\
    If $n$ is even let right be a heads and let up be a tails if we want there to be an equal number of heads and tails then we want $n/2$ heads and $n/2$ tails thus this counds the number of lattice paths from $0,0$ to $n/2,n/2$
    thus this is $\binom{n}{n/2}$ thus\\
    $\P(\# \t{ of heads and tails are equal on $n$-th throw})=\dfrac{\binom{n}{n/2}}{2^n}$
    \item $\P(\t{exactly two heads have appeared to date})$\\
    Total number of posibilities is $2^n$\\
    The number of ways we get exaclty two heads is just $\binom{n}{2}$ since we are picking 2 heads and the rest are tails thus\\
    $\P(\t{exactly two heads have appeared to date})=\dfrac{\binom{n}{2}}{2^n}$
    \item $\P(\t{at least two heads have appeared to date})$\\
    Total number of posibilities is $2^n$\\
    The number of ways to get at least two heads is also to count the complement of the ways to get 0 heads and 1 head thus\\
    the number of ways to get 1 head or 0 heads is $1+n$ thus $\P(\t{at least two heads})=1-\P(\t{1 or 0 heads})=1-\dfrac{n+1}{2^n}$
\end{enumerate}
\end{document}